\section{Algorytmy genetyczne}

Algorytmy genetyczne (ang. \textit{genetic algorithms} --- GA) to najszerzej znana gałąź algorytmów ewolucyjnych. Tradycyjnie operują one na binarnych łańcuchach genów, ale współczesne warianty stosują również reprezentacje liczb rzeczywistych, permutacyjne lub inne, dostosowane do struktury danego problemu \cite{Katoch2021}, a kluczowymi operatorami są: \textbf{krzyżowanie}, czyli wymiana fragmentów rozwiązań między dwoma rodzicami oraz \textbf{mutacja}, czyli losowe modyfikacje genomu potomka z małym prawdopodobieństwem. Selekcja odbywa się metodą ruletki, turniejową lub rankingową, faworyzując osobniki o wyższym przystosowaniu. Analiza Katocha i in. \cite{Katoch2021} pokazuje, że mimo upływu dekad od pierwszego sformułowania algorytmu, GA nadal jest aktywnym obszarem badań, z licznymi modyfikacjami poprawiającymi równowagę między eksploracją a eksploatacją oraz odporność na przedwczesną zbieżność.

Algorytmy genetyczne znalazły również zastosowanie w kontekście gier karcianych. Yang i in. \cite{yang2021} zastosowali GA do optymalizacji składu talii w grze Legends of Code and Magic, kodując talię jako binarny wektor obecności kart i osiągając wyniki lepsze niż konstruowanie talii w sposób losowy. Podejście to można przenieść na grunt Hearthstone, gdzie problem budowania talii ma podobną strukturę kombinatoryczną \cite{kowalski2023}. Algorytmy genetyczne są jednak rzadziej stosowane do optymalizacji wag ciągłej funkcji oceny stanu planszy, gdzie lepiej sprawdzają się algorytmy operujące bezpośrednio na przestrzeniach liczb rzeczywistych \cite{sanchez2020}.
